\section{Các trường hợp thu hồi được}
\label{sec:ch06-thu-hoi}

\index{mạng ReLU}%
\index{partial dependence plot}%
\tn{Độ đo}{measure} làm rõ giả định, và ba trường hợp của mục trước cho thấy
khung không đứng tách khỏi các công cụ quen thuộc: attribution dựa trên
\tn{loại bỏ đặc trưng}{feature removal}, biến thể giả định độc lập và partial
dependence plot đều viết được như kỳ vọng của đầu ra mô hình dưới một cách lấy
trung bình đã chọn~\autocite{p17firstprinciples}.

Với hàm từng đoạn affine, bài báo còn đưa attribution về tổng trên các vùng
affine. Mạng ReLU sâu thuộc lớp hàm này: trên mỗi vùng kích hoạt cố định,
mạng là một hàm affine; các hệ số của vùng có thể biểu diễn từ trọng số mạng.
Vì vậy, dạng tổng quát thu về một dạng đóng cho mạng ReLU sau khi phân hoạch
không gian đầu vào theo các vùng đó: mỗi vùng $P$ đóng góp trọng lượng
$\mu_{j,x}(P)$ nhân với giá trị của hàm affine của vùng ấy tại trọng tâm của
độ đo trên vùng, trong đó trọng tâm là điểm trung bình của $y$ theo
$\mu_{j,x}$ giới hạn trên $P$~\autocite{p17firstprinciples}. Hệ quả này hẹp hơn phương
trình~\ref{eq:ch06-attribution-do-do} một bậc: nó đòi hỏi mỗi độ đo $\mu_{j,x}$
là độ đo dương, tức chỉ gán trọng lượng không âm cho mọi vùng, thay vì độ đo có
dấu, vốn được phép gán trọng lượng âm cho một số vùng. Ở mức
\tn{quy tắc nguyên tử}{atomic attribution rule}, điều đó nghĩa là attribution
của \tn{hàm chỉ báo}{indicator function} không bao giờ âm. Đổi lại, bài báo
ghi rằng với attribution dương thì điều kiện ổn định trước hội tụ đều của
mục~\ref{sec:ch06-tu-nguyen-tu-den-tich-phan} tự thỏa~\autocite{p17firstprinciples}.

Từ đây không nên suy ra rằng mọi attribution thực hành đều đã có một thuật
toán rẻ. Bài báo nêu hai cách xấp xỉ tích phân: Monte Carlo với một bộ lấy mẫu, dành
cho mô hình có số chiều lớn, hoặc tổng Riemann cổ điển được vector hóa trên
GPU. Cả hai vẫn phụ thuộc vào độ đo
đã chọn, số mẫu hoặc độ mịn của phân hoạch. Khung cho phép so những quyết định
ấy trong cùng một bộ ký hiệu; nó không làm chi phí tính toán và sai lệch xấp xỉ
biến mất.
